\section{Results}\label{sec:results}

\subsection{Frozen features on the pretraining classes}\label{sec:frozen}

Coarser labels leave the frozen features less able to tell the pretraining classes apart (table~\ref{tab:probes} and
figure~\ref{fig:accuracy}).
A linear probe on the features of the \Kfull-class model identifies the \Kfull{} classes with an accuracy of
\VLpJcTwoFull\%, as does the checkpoint released with Sophon (\VLpJcTwoSophon\%); chance is \VChance\% and an untrained
network gives \VLpJcTwoUntrained\%.
Merging the QCD classes changes little: the \Kres-class model reaches \VLpJcTwoRes\%, and its AUC is lower only over the
QCD classes it merges (table~\ref{tab:pairs}).
The \Kmid- and \Kcoarse-class sets, which merge flavors within the resonance decays, fall to \VLpJcTwoMid\% and
\VLpJcTwoCoarse\%, a loss of \VDLpJcTwoCoarse{} percentage points in all.
These accuracies are far above what the coarse labels alone allow.
A classifier that knows the \Kcoarse-class label of a jet and picks uniformly among the classes it contains would reach
\VLabelOnlyCoarse\%, and one that knows the \Kmid-class label \VLabelOnlyMid\%, so the coarse features keep much of the
fine information that their labels discard.
\vfig{fig2_accuracy}{Accuracy of the nested, random and flavor-boundary label sets on the \Kfull{} JetClass-II classes
(a) and on JetClass (b). The upper parts give linear probes on the frozen features (open markers), the lower parts
models fine-tuned on \FtSizeESix{} jets (filled), each part with its own range. The flavor-boundary groupings
(section~\ref{sec:labelsets}) have no flavor boundary, a $b$-quark boundary or a random one. Markers and bars are means
and standard deviations over runs, faint dots the single runs and, for the random groupings, the \VRandWord{} draws,
each a single run. The dashed line and band are the network trained from scratch.}{fig:accuracy}
\input{tables/table2}

The loss falls mainly on the classes that a label set merges (figure~\ref{fig:classes}(a) and
figure~\ref{fig:recall}).
At \Kmid{} classes the mean recall changes by $\VRecallMidQcd$ points for the QCD classes, $\VRecallMidQuark$ for
classes merged only across quark flavor, $\VRecallMidLepton$ for classes merged only across lepton flavor,
$\VRecallMidBoth$ for classes merged across both, and $\VRecallMidNeither$ for classes the set leaves unmerged.
At \Kcoarse{} classes, where the remaining quark flavor is merged as well, the same changes are $\VRecallCoarseQcd$,
$\VRecallCoarseQuark$, $\VRecallCoarseLepton$, $\VRecallCoarseBoth$ and $\VRecallCoarseNeither$.
The QCD classes carry the same label in the \Kres-, \Kmid- and \Kcoarse-class sets, yet lose more recall as the resonant
labels coarsen, which suggests that part of the loss reaches classes whose own labels do not change.

A probe of the \Kfull{} classes might seem to favor fine labels by construction.
The reverse test does not favor coarse ones: a linear probe asked only for the \Kcoarse{} groups of the coarsest set is
as accurate on the features of the \Kfull-class model as on those of the \Kcoarse-class model itself (\VGroupsLpFull\%
and \VGroupsLpCoarse\%), and an MLP probe is slightly more accurate on the \Kfull-class features (\VGroupsMlpFull\%
against \VGroupsMlpCoarse\%; table~\ref{tab:groups}).

\subsection{Pairs of classes}\label{sec:pairs}

Merging two classes raises the error of a probe that separates them, up to fourfold for bottom against charm
(figure~\ref{fig:classes}(b), table~\ref{tab:pairs}; rejection curves in figures~\ref{fig:otherpairs}
and~\ref{fig:qcdpairs}).
The decays $X\to b\bar b$ and $X\to c\bar c$ differ only in the flavor of their quarks and share a class only in the
\Kcoarse-class set.
A linear probe separates them with an AUC of \VAucBbFull{} on the features of the \Kfull-class model and
\VAucBbCoarse{} on those of the \Kcoarse-class model, and the four-quark decays $X\to b\bar bqq$ and $X\to c\bar cqq$
behave alike, with \VAucBbqqFull{} against \VAucBbqqCoarse.
The decay $X\to bc$ shares a class with $X\to bq$ from \Kmid{} classes on, and its separation from $X\to bq$ and
$X\to cs$ falls from \VAucBcFullThree{} at \Kfull{} classes to \VAucBcMidThree{} at \Kmid{} and \VAucBcCoarseThree{} at
\Kcoarse.
For $X\to gg$ against $X\to q\bar q$, merged from \Kmid{} classes on, the change is small: the AUC falls from
\VAucGgFull{} to \VAucGgCoarse{} at \Kcoarse{} classes, an error \VErrGgCoarsePct\% larger.
\vfig{fig3_classes}{Where the frozen-feature accuracy is lost. (a) Mean change of per-class recall under a linear probe
against the \Kfull-class model, by what the label set merges a class with; bars are standard deviations over runs.
(b) One minus the AUC of a binary linear probe for pairs of classes, relative to the \Kfull-class model (value in
parentheses); open where the labels merge the pair, half-filled where they merge part of it. The pair $X\to gg$ against
$X\to q\bar q$ has fewer test jets (appendix~\ref{app:training}) and a larger uncertainty.}{fig:classes}

A distinction that the labels keep is not fully protected.
At \Kmid{} classes $X\to b\bar b$ and $X\to c\bar c$ are still separate classes, yet the error $1-\mathrm{AUC}$ is
\VErrBbMidPct\% larger than with the \Kfull-class model (\VErrBbMid{} against \VErrBbFull), and \VErrBbqqMidPct\%
larger for $X\to b\bar bqq$ against $X\to c\bar cqq$.

\subsection{Which classes are merged}\label{sec:which}

Label sets of the same size differ widely.
The random groupings have as many classes as the \Kcoarse-class set, of the same sizes, yet a linear probe on their
features reaches \VLpJcTwoRand\%, above the \Kmid-class model and \VDLpJcTwoRand{} points below the \Kfull-class model
(figure~\ref{fig:accuracy}).
Merging a pair inside a random grouping hardly affects it: the groupings that place $X\to b\bar b$ and $X\to c\bar c$ in
one class separate them with an AUC of \VAucBbRandMerged, against \VAucBbRandApart{} for the groupings that keep them
apart and \VAucBbCoarse{} for the \Kcoarse-class model.
In accuracy the flavor-boundary groupings lie between the \Kcoarse-class set and the random groupings
(table~\ref{tab:flavor}).
The grouping with no flavor boundary reaches \VLpJcTwoFlavZero\%, \VDLpJcTwoFlavZero{} points above the \Kcoarse-class
set, and moving one set of four-quark decays gives \VLpJcTwoFlavOne\% for the $b$-quark boundary and \VLpJcTwoFlavOneR\%
for a random one.

The frozen accuracy rises with the separation $S$ of equation~(\ref{eq:separation}) across all \VSepNWord{} label sets
(figure~\ref{fig:separation}).
$S$ runs from \VSepCoarse{} for the \Kcoarse-class set through \VSepFlavZero--\VSepFlavOneR{} for the flavor-boundary
groupings and \VSepRandLo--\VSepRandHi{} for the random groupings to \VSepRes{} at \Kres{} classes, with a Spearman
correlation of \RevSepRho{} ($p=\RevSepP$) between $S$ and the accuracy.
Weighting the class pairs by the confusion of an untrained network, which shares no pretraining with any label set,
gives a correlation of \RevSepUntrainedRho, and counting separated pairs without weights gives \RevSepNaiveRho.
The random groupings cut across flavor families, so they keep apart most of the confusable pairs that the
\Kcoarse-class set merges.
The correlation is not perfect.
The \Kmid-class set keeps more of the weighted separation (\VSepMid) than the groupings with no flavor boundary and with a
$b$-quark boundary (both \VSepFlavZero), yet its accuracy is lower (\VLpJcTwoMid\% against \VLpJcTwoFlavZero\% and
\VLpJcTwoFlavOne\%), and among the random groupings, pretrained once each, the accuracy does not follow the spread
in~$S$.
\vfig{fig_separation}{Frozen-feature accuracy rises with the separation a label set keeps between confusable classes.
Each point is a pretrained label set: its separation $S$ (equation~(\ref{eq:separation})) against the linear-probe
accuracy on the \Kfull{} JetClass-II classes, the mean and standard deviation over runs; each random grouping is a
single run. Over the \VSepNWord{} label sets the Spearman correlation is $\rho=\RevSepRho$ ($p=\RevSepP$); weighting by
the confusion of an untrained network instead gives $\rho=\RevSepUntrainedRho$.}{fig:separation}

\subsection{Resonance benchmark}\label{sec:resonance}

Coarse labels also lower the signal efficiency in Sophon's resonance benchmark (figure~\ref{fig:resonance} and
table~\ref{tab:resonance}).
At a background efficiency of \VXbxEpsB\%, a classifier trained on the frozen features of the \Kfull-class model
selects $X\to b\bar b$ with a signal efficiency of \VXbbFull\% and $X\to bs$ with \VXbsFull\%.
With the \Kcoarse-class model these fall to \VXbbCoarse\% and \VXbsCoarse\%, losses of \RevDXbbCoarse{} and
\RevDXbsCoarse{} points, and the \Kmid-class model lies in between, at \VXbbMid\% and \VXbsMid\%.
The loss is larger for $X\to bs$, which is a class in none of the label sets.
For $X\to b\bar b$ the two finest label sets also allow Sophon's own selection, the sum of equation~(\ref{eq:sum}) over
the signal classes, which reaches \VXbbDirectFullOne\% with the \Kfull-class model against \VXbbFull\% for the trained
classifier.
Sophon's network trained from scratch on the same jets reaches \VXbbScratch\% for $X\to b\bar b$, with runs that differ
strongly at low background efficiency, and \VXbsScratch\% for $X\to bs$; the \Kcoarse-class features lie
\RevDXbbScratch{} and \RevDXbsScratch{} points above it, within the uncertainties.
\vfig{fig5_resonance}{Sophon's resonance benchmark for $X\to b\bar b$ (a) and $X\to bs$ (b): QCD rejection against
signal efficiency in the mass window. Solid lines are a classifier trained on the frozen features of each label set
and, in grey, Sophon's network trained from scratch on the same jets; dash-dotted lines use the output scores where
the label set keeps $X\to b\bar b$. Bands are standard deviations over runs; the dotted line marks \VXbxEpsB\%
background efficiency. Lower panels give the ratio to the \Kfull-class model.}{fig:resonance}

\subsection{Fine-tuning}\label{sec:finetune}

Fine-tuning narrows the gap between label sets but does not close it (table~\ref{tab:finetuned} and
figure~\ref{fig:accuracy}).
On JetClass-II the frozen \Kfull-class features are the natural ceiling, since they were trained on these very labels.
After fine-tuning on \FtSizeESix{} jets the \Kcoarse-class model reaches \VFtJcTwoCoarse\%, the \Kmid-class model
\VFtJcTwoMid\% and the \Kres-class model \VFtJcTwoRes\%, against \VFtJcTwoFull\% for the \Kfull-class model and
\VFtJcTwoScratch\% for the same network trained from scratch.
The \VDLpJcTwoCoarse-point frozen gap between the finest and the coarsest labels shrinks to \VDFtJcTwoCoarse{} points,
which is still more than a third of the \VDFtJcTwoGain{} points that pretraining on \Kfull{} classes gains over training
from scratch.
With our recipe, which replaces the output layer with a new one trained at \MlstFtHeadMult{} times the backbone's
learning rate, the fine-tuned \Kfull-class model ends \VDCsLpGap{} points below its own linear probe
(appendix~\ref{app:training}).

On JetClass, whose ten classes are not pretraining labels and which was simulated without pileup, the differences are
smaller.
With frozen features the \Kfull-class model reaches \VLpJcOneFull\% and the \Kcoarse-class model \VLpJcOneCoarse\%,
\RevDLpJcOneCoarse{} points apart; after fine-tuning they reach \VFtJcOneFull\% and \VFtJcOneCoarse\%,
\RevDFtJcOneCoarse{} points apart, while pretraining itself gains \VDFtJcOneGain.
Per class, the frozen loss is largest for the decays of the Higgs boson to charm and to bottom quarks, \VDLpRecallHcc{}
and \VDLpRecallHbb{} points of recall, the two JetClass classes separated only by quark flavor.
After fine-tuning the loss is smaller and spread over several classes, and only the decay to bottom quarks, at
\VDFtRecallHbb{} points, is resolved on its own (table~\ref{tab:perclass_jc1}).
The random groupings, which kept most of the accuracy on the pretraining classes, reach \VLpJcOneRand\% on JetClass
with frozen features, \VDLpJcOneRand{} points below the \Kfull-class model and \RevDLpJcOneRandCoarse{} points below the
\Kcoarse-class model.
After fine-tuning they reach \VFtJcTwoRand\% on JetClass-II, level with the \Kmid-class set, and \VFtJcOneRand\% on
JetClass, between the \Kmid- and \Kcoarse-class sets.

The label set matters more when labeled jets are scarce (figure~\ref{fig:size}).
On JetClass-II the \Kcoarse-class model trails the \Kfull-class model by \VDFtJcTwoCoarseSmallest{} points at
\FtSizeEFour{} training jets, \VDFtJcTwoCoarseSmaller{} at \FtSizeEFive{} and \VDFtJcTwoCoarse{} at \FtSizeESix, and
linear probes trained on fewer jets keep the order of the nested sets (figure~\ref{fig:probesize}).
\vfig{fig4_training_size}{Fine-tuned accuracy on JetClass-II (a) and JetClass (b) against the number of training jets,
for each label set and for the network trained from scratch. Lines and bands are means and standard deviations over
runs (values in table~\ref{tab:size}). Lower panels give each label set's difference to the \Kfull-class model in
percentage points, with both run spreads combined.}{fig:size}
\input{tables/table3}

\subsection{Binary benchmarks}\label{sec:binary}

On top tagging and quark--gluon discrimination pretraining itself adds little after fine-tuning on the full training
sets: the \Kfull-class model reaches \VFtTopAccFull\% on top tagging against \VFtTopAccScratch\% for the network
trained from scratch, and \VFtQgAccFull\% against \VFtQgAccScratch\% on quark--gluon discrimination
(table~\ref{tab:finetuned}).
Differences between label sets are bounded by this small total, and the four nested sets agree to within
\VFtTopAccRange{} points in top-tagging accuracy and \VFtQgAccRange{} points in quark--gluon accuracy, where relative
differences above \RevReachTopFtAccLo--\RevReachTopFtAccHi\% and \RevReachQgFtAccLo--\RevReachQgFtAccHi\% of the
accuracy would stand out from the run spread.
The background rejection is less precise: a difference would have to exceed
\RevReachTopFtRthirtyLo--\RevReachTopFtRthirtyHi\% of $\mathrm{Rej}_{30\%}$ on top tagging and
\RevReachQgFtRthirtyLo--\RevReachQgFtRthirtyHi\% on quark--gluon discrimination.
The top-tagging rejection of the \Kfull-class model rises from \VFtTopRejScratch{} to \VFtTopRejFull{} over training
from scratch (difference \RevDFtTopRejGain), with about a hundred background jets passing at this working point.
With frozen features (table~\ref{tab:probes}) the nested sets lie within \RevSpreadTopLpAccPts{} points of each other in
top-tagging accuracy and \RevSpreadQgLpAccPts{} points in quark--gluon accuracy, with the \Kcoarse-class set highest on
top tagging, and fine-tuning on fewer jets shows no ordering by label set either (figure~\ref{fig:ftbinary}).
The random groupings lie below every nested set on frozen top tagging, a difference resolved only against the
\Kcoarse-class set.
Neither data set carries the track displacements that identify heavy flavor (section~\ref{sec:tasks}).
Quark--gluon discrimination does draw on a distinction that the coarser sets merge, gluons against light quarks, but that
merge costs little even on the pretraining classes (section~\ref{sec:pairs}), and the benchmark shows no dependence on
the label set, also on jets simulated with Herwig~\cite{herwigsevenone} instead of Pythia~\cite{pythia}
(table~\ref{tab:herwig}).

\subsection{Jet mass}\label{sec:mass}

The auxiliary regression makes the jet mass much easier to recover from the frozen features (table~\ref{tab:mass}).
An MLP probe on the features of the \Kres-class model estimates the mass with a resolution of \VResMlpRes\%, against
\VResSoftDrop\% for the reconstructed soft-drop mass; with the regression in pretraining the same probe reaches
\VResMlpResMass\%, and the regression output of the network \VResHeadResMass\%.
At \Kcoarse{} classes the probe improves from \VResMlpCoarse\% to \VResMlpCoarseMass\%, and the regression output is
sharper than the soft-drop mass over the whole range of jet $p_{\mathrm T}$ and mass (figure~\ref{fig:masspt}).
Without the regression the \Kcoarse-class model already gives a better resolution than the \Kfull-class model,
\VResMlpCoarse\% against \VResMlpFull\%, which suggests that labels organized by decay structure reward features tied
to the jet mass.
\input{tables/table4}

The price in classification is small.
With frozen features the regression costs \RevDLpMassRes{} points on the pretraining classes at \Kres{} classes; at
\Kcoarse{} classes the two runs with the regression differ by several points (\VLpJcTwoCoarseMassSOne\% and
\VLpJcTwoCoarseMassSTwo\%), too much to measure a cost.
After fine-tuning the models with the regression are lower on JetClass-II by \RevDFtMassRes{} points at \Kres{}
classes and \RevDFtMassCoarse{} points at \Kcoarse{} classes, and on JetClass by about \VDFtJcOneMassResOne{} points
(\RevDFtJcOneMassCoarse{} at \Kcoarse{} classes).
Because top tagging relies on the jet mass, the regression also raises the frozen-feature rejection there, from
\VLpTopRejCoarse{} to \VLpTopRejCoarseMass{} at \Kcoarse{} classes and from \VLpTopRejRes{} to \VLpTopRejResMass{} at
\Kres{} (table~\ref{tab:probes}), with the MLP probe in the same direction (table~\ref{tab:mlp}); neither difference is
resolved.
